\chapter{Bẫy Bạn Bè Và So Sánh}
\markboth{Bẫy Bạn Bè Và So Sánh}{The Trap of Friends and Comparison}

\section{Nhìn Thành Tựu Người Khác, Quên Nhịp Của Mình}
\begin{truyen}{Nhìn Thành Tựu Người Khác, Quên Nhịp Của Mình}{Watching Others and Forgetting Your Own Pace}
\chuhoa{C}{ó người trẻ mở điện thoại ra là thấy ai đó được nhận việc, ai đó mua xe, ai đó cưới, ai đó đi xa hơn mình.}
\textit{(Some young people open their phone and immediately see someone getting hired, buying a car, marrying, or moving further ahead.)}

Những hình ảnh đó dần làm họ quên mất đời mình có nhịp riêng.
\textit{(Those images slowly make them forget that life has individual timing.)}

Khi so sánh trở thành thói quen, mỗi bước của mình dù vẫn tốt cũng thấy nhạt đi.
\textit{(When comparison becomes a habit, even your own solid progress starts to feel dull.)}

Người luôn nhìn ngang để đo giá trị mình rất khó đi thẳng lâu dài.
\textit{(People who constantly look sideways to measure their worth struggle to keep walking straight.)}

\end{truyen}

\begin{baihoc}
Không phải cứ chậm hơn người khác là sai nhịp. Có khi em chỉ đang đi đúng đời mình.
\textit{(Being slower than others does not mean you are off track. You may simply be walking your own life.)}
\end{baihoc}

\begin{ghinhoanh}
\textbf{Short English to remember:}

Being slower than others does not mean you are off track.

You may simply be walking your own life.

\end{ghinhoanh}

\begin{tuhocnhanh}
\textbf{Câu hỏi tự học / Self-study questions}

1. Nhân vật đang quên điều gì? \textit{What is the character forgetting?}

2. So sánh làm mờ điều gì trong mình? \textit{What does comparison blur inside you?}

3. Nhịp riêng của em đang là gì? \textit{What is your own pace right now?}
\end{tuhocnhanh}
\ngancach

\section{Chơi Với Nhóm Chỉ Biết Đua Hình Ảnh}
\begin{truyen}{Chơi Với Nhóm Chỉ Biết Đua Hình Ảnh}{Staying in a Group That Competes Through Image}
\chuhoa{C}{ó những nhóm bạn gặp nhau là so đồ mới, chỗ check-in, ai được chú ý hơn, ai có vẻ sống tốt hơn.}
\textit{(Some friend groups gather mainly to compare new things, check-in spots, attention, and who seems to be living better.)}

Lúc đầu không khí đó vui và kích thích.
\textit{(At first, that atmosphere feels exciting.)}

Nhưng lâu dần, nó khiến người ta tiêu nhiều hơn, diễn nhiều hơn, và mệt nhiều hơn chỉ để không tụt hạng vô hình.
\textit{(Over time, it makes people spend more, perform more, and feel more tired just to avoid dropping in an invisible ranking.)}

Một nhóm bạn không cho em bình yên làm chính mình thì cái vui trong nhóm đó thường rất đắt.
\textit{(If a friend group does not allow you to be peaceful as yourself, the fun inside it is usually expensive.)}

\end{truyen}

\begin{baihoc}
Bạn bè tốt không bắt em phải diễn để được ở lại.
\textit{(Good friends do not require a performance for you to belong.)}
\end{baihoc}

\begin{ghinhoanh}
\textbf{Short English to remember:}

Good friends do not require a performance for you to belong.

\end{ghinhoanh}

\begin{tuhocnhanh}
\textbf{Câu hỏi tự học / Self-study questions}

1. Nhóm bạn này đang nuôi điều gì? \textit{What is this group feeding?}

2. Vì sao cái vui này trở nên đắt? \textit{Why does this fun become expensive?}

3. Em có đang phải diễn để hòa nhập không? \textit{Are you performing just to fit in?}
\end{tuhocnhanh}
\ngancach

\section{Ngại Thành Công Vì Sợ Lạc Nhóm}
\begin{truyen}{Ngại Thành Công Vì Sợ Lạc Nhóm}{Holding Back Success to Avoid Outgrowing the Group}
\chuhoa{C}{ó người thật ra làm được nhiều hơn, nhưng lại tự kìm vì sợ nếu tiến nhanh quá sẽ khác bạn bè mình.}
\textit{(Some people could actually go much further, but they hold themselves back for fear of becoming too different from their friends.)}

Nỗi sợ lạc nhóm khiến họ tự giảm chuẩn hoặc làm lơ cơ hội.
\textit{(The fear of no longer belonging makes them lower their own standard or ignore opportunities.)}

Tình bạn đẹp không phải là cùng đứng yên cho đỡ lệch.
\textit{(Healthy friendship is not staying equally still so no one feels different.)}

Nếu một bước lớn lên của em đe dọa được một tình bạn, có thể tình bạn đó đang yếu hơn em nghĩ.
\textit{(If your growth threatens a friendship, that friendship may be weaker than you thought.)}

\end{truyen}

\begin{baihoc}
Đừng tự làm nhỏ mình chỉ để vẫn vừa trong một vòng tròn cũ.
\textit{(Do not shrink yourself just to keep fitting inside an old circle.)}
\end{baihoc}

\begin{ghinhoanh}
\textbf{Short English to remember:}

Do not shrink yourself just to keep fitting inside an old circle.

\end{ghinhoanh}

\begin{tuhocnhanh}
\textbf{Câu hỏi tự học / Self-study questions}

1. Nhân vật đang kìm điều gì? \textit{What is the character holding back?}

2. Nỗi sợ nào đứng phía sau? \textit{What fear stands behind it?}

3. Em có đang làm nhỏ mình cho vừa với ai không? \textit{Are you shrinking yourself to fit someone else?}
\end{tuhocnhanh}
\ngancach

\section{Lấy Lời Khen Trong Nhóm Làm Thước Đo}
\begin{truyen}{Lấy Lời Khen Trong Nhóm Làm Thước Đo}{Using Group Approval as the Main Measure}
\chuhoa{N}{hiều người chỉ thấy mình ổn khi nhóm khen, nhóm công nhận, hoặc nhóm còn chú ý tới mình.}
\textit{(Many people only feel okay when the group praises them, validates them, or keeps paying attention.)}

Khi lời khen đó mất đi, họ rơi rất nhanh vào nghi ngờ bản thân.
\textit{(When that approval disappears, they fall quickly into self-doubt.)}

Thước đo đặt bên ngoài làm giá trị bản thân lên xuống theo ánh mắt người khác.
\textit{(A measure placed outside yourself makes your self-worth rise and fall with other people’s eyes.)}

Người trưởng thành cần một chỗ đứng bên trong đủ chắc để không phải xin phép đám đông mới thấy mình có giá trị.
\textit{(A mature person needs an inner footing strong enough that worth does not depend on crowd permission.)}

\end{truyen}

\begin{baihoc}
Lời khen có thể vui, nhưng nó không nên là móng nhà của lòng tự trọng.
\textit{(Praise can feel good, but it should not become the foundation of self-respect.)}
\end{baihoc}

\begin{ghinhoanh}
\textbf{Short English to remember:}

Praise should not become the foundation of self-respect.

\end{ghinhoanh}

\begin{tuhocnhanh}
\textbf{Câu hỏi tự học / Self-study questions}

1. Nhân vật đang đo mình bằng gì? \textit{What is the character using as a measure?}

2. Vì sao thước đo đó bấp bênh? \textit{Why is that measure unstable?}

3. Em cần xây chỗ đứng bên trong bằng điều gì? \textit{What can help you build stronger inner footing?}
\end{tuhocnhanh}
\ngancach

\section{Thấy Bạn Đi Nhanh, Tưởng Mình Đi Sai}
\begin{truyen}{Thấy Bạn Đi Nhanh, Tưởng Mình Đi Sai}{Seeing Friends Move Fast and Assuming You Are on the Wrong Road}
\chuhoa{C}{ó người chỉ cần nhìn bạn bè có việc tốt, thu nhập tốt, hay nhiều thành tựu sớm là bắt đầu thấy nhịp của mình có vấn đề.}
\textit{(Some people only need to see friends gain jobs, income, or early achievements to start doubting their own pace.)}

Họ quên rằng mỗi người đang khởi hành từ một điểm khác nhau và mang một hoàn cảnh khác nhau.
\textit{(They forget that each person begins from a different point and carries a different set of circumstances.)}

So sánh tốc độ mà không nhìn đường đi rất dễ biến cố gắng của mình thành một cảm giác thua thiệt oan.
\textit{(Comparing speed without seeing the path can unfairly turn your own effort into a feeling of loss.)}

Đi chậm chưa chắc là đi sai. Có khi em chỉ đang đi một con đường cần nền chắc hơn.
\textit{(Moving slowly does not mean moving wrongly. Sometimes you are simply on a path that needs a stronger foundation.)}

\end{truyen}

\begin{baihoc}
Đừng lấy tốc độ của người khác làm bản án cho nhịp của mình.
\textit{(Do not use someone else's speed as a verdict against your own pace.)}
\end{baihoc}

\begin{ghinhoanh}
\textbf{Short English to remember:}

Someone else's speed is not a verdict on your pace.

\end{ghinhoanh}

\begin{tuhocnhanh}
\textbf{Câu hỏi tự học / Self-study questions}

1. Nhân vật đang quên điều gì khi so sánh? \textit{What is the character forgetting while comparing?}

2. Vì sao đi chậm chưa chắc là sai? \textit{Why does moving slowly not necessarily mean being wrong?}

3. Em có đang kết án nhịp của mình vì nhìn nhịp người khác không? \textit{Are you judging your pace by someone else's?}
\end{tuhocnhanh}
\ngancach

\section{Biến Cuộc Gặp Gỡ Thành Cuộc Thi Ngầm}
\begin{truyen}{Biến Cuộc Gặp Gỡ Thành Cuộc Thi Ngầm}{Turning Every Gathering into a Silent Contest}
\chuhoa{C}{ó người đi gặp bạn bè mà lòng không thật sự thư giãn, vì lúc nào cũng âm thầm so điểm: ai giỏi hơn, ai hơn hình ảnh, ai hơn thành tích.}
\textit{(Some people meet friends without really relaxing because they are silently scoring everyone: who looks better, earns more, or appears ahead.)}

Cuộc gặp đáng ra để nối lại tình thân lại biến thành một bài kiểm tra tự ái.
\textit{(A gathering meant for connection turns into an ego exam.)}

Khi tình bạn bị kéo vào logic thắng thua, niềm vui gặp nhau sẽ bị hao đi rất nhanh.
\textit{(When friendship is pulled into the logic of winning and losing, the joy of meeting fades quickly.)}

Không phải cuộc gặp nào cũng cần một người nổi bật nhất.
\textit{(Not every meeting needs one person to stand out the most.)}

\end{truyen}

\begin{baihoc}
Tình bạn khỏe không cần bảng xếp hạng ngầm.
\textit{(Healthy friendship does not need a hidden ranking board.)}
\end{baihoc}

\begin{ghinhoanh}
\textbf{Short English to remember:}

Healthy friendship does not need hidden rankings.

\end{ghinhoanh}

\begin{tuhocnhanh}
\textbf{Câu hỏi tự học / Self-study questions}

1. Cuộc gặp đang bị biến thành điều gì? \textit{What is the gathering being turned into?}

2. Điều gì hao đi nhanh nhất? \textit{What disappears first?}

3. Em có đang âm thầm chấm điểm bạn bè không? \textit{Do you silently score your friends?}
\end{tuhocnhanh}
\ngancach

\section{Ganh Tị Nhưng Giả Vờ Bình Thản}
\begin{truyen}{Ganh Tị Nhưng Giả Vờ Bình Thản}{Feeling Envy While Pretending to Be Fine}
\chuhoa{C}{ó người nghe tin vui của bạn rồi mỉm cười rất đẹp, nhưng trong lòng nặng xuống mà không dám gọi tên.}
\textit{(Some people hear a friend's good news and smile gracefully while feeling their heart sink inside.)}

Vì ngại xấu hổ, họ không thừa nhận đó là ghen tị mà chỉ để cảm giác ấy âm thầm tích lại.
\textit{(Because of shame, they do not admit it is envy and let the feeling quietly accumulate.)}

Điều không được gọi tên thường khó được chữa lành.
\textit{(What is not named is hard to heal.)}

Ghen tị được nhìn thẳng có thể thành động lực sửa mình. Ghen tị bị giấu lâu rất dễ thành cay đắng.
\textit{(Envy faced honestly can become motivation for growth. Envy hidden for too long turns bitter.)}

\end{truyen}

\begin{baihoc}
Thừa nhận một cảm xúc xấu không làm em xấu hơn. Nó cho em cơ hội xử lý nó.
\textit{(Admitting an ugly emotion does not make you uglier. It gives you a chance to process it.)}
\end{baihoc}

\begin{ghinhoanh}
\textbf{Short English to remember:}

Named envy can become growth; hidden envy becomes bitterness.

\end{ghinhoanh}

\begin{tuhocnhanh}
\textbf{Câu hỏi tự học / Self-study questions}

1. Nhân vật đang giấu cảm xúc nào? \textit{What feeling is the character hiding?}

2. Vì sao cần gọi tên nó? \textit{Why does it need to be named?}

3. Em thường làm gì với cảm giác ganh tị? \textit{What do you usually do with envy?}
\end{tuhocnhanh}
\ngancach

\section{Chọn Bạn Theo Độ Hào Nhoáng}
\begin{truyen}{Chọn Bạn Theo Độ Hào Nhoáng}{Choosing Friends by Their Shine}
\chuhoa{C}{ó người thích đứng gần những nhóm nhìn rất nổi bật vì cảm giác như mình cũng sáng lên theo.}
\textit{(Some people like staying near shiny, impressive groups because it makes them feel brighter too.)}

Họ ít hỏi nhóm đó có tử tế, có chân thành, hay có giúp nhau lớn lên không.
\textit{(They rarely ask whether that group is kind, sincere, or truly helps each other grow.)}

Hào nhoáng tạo ra cảm giác có vị thế rất nhanh, nhưng không chắc tạo ra chỗ dựa thật.
\textit{(Shine creates the feeling of status quickly, but not necessarily real support.)}

Bạn bè đáng giữ không phải là người làm em nhìn sang hơn, mà là người giúp em sống tốt hơn.
\textit{(Friends worth keeping are not those who make you look better, but those who help you live better.)}

\end{truyen}

\begin{baihoc}
Ánh sáng bên ngoài của một nhóm không thay được chất lượng bên trong của tình bạn.
\textit{(A group's external shine cannot replace the inner quality of friendship.)}
\end{baihoc}

\begin{ghinhoanh}
\textbf{Short English to remember:}

Shine cannot replace the inner quality of friendship.

\end{ghinhoanh}

\begin{tuhocnhanh}
\textbf{Câu hỏi tự học / Self-study questions}

1. Nhân vật đang bị hút bởi điều gì? \textit{What is the character attracted to?}

2. Điều gì chưa được hỏi tới? \textit{What has not been examined?}

3. Em chọn bạn vì hào nhoáng hay vì chất lượng sống? \textit{Do you choose friends for shine or for quality of life?}
\end{tuhocnhanh}
\ngancach

\section{Để Một Nhóm Quyết Định Mức Mơ Của Mình}
\begin{truyen}{Để Một Nhóm Quyết Định Mức Mơ Của Mình}{Letting a Group Decide the Size of Your Dream}
\chuhoa{C}{ó người từ khi chơi trong một nhóm nào đó thì dần thu nhỏ hoặc phóng to ước mơ của mình chỉ để không lệch khỏi mặt bằng chung.}
\textit{(Some people slowly shrink or enlarge their dreams just to stay aligned with the average of a group.)}

Họ sợ mơ lớn quá sẽ bị nói là viển vông, còn mơ khác quá sẽ bị lạc nhịp.
\textit{(They fear that dreaming too big will seem unrealistic, and dreaming too differently will make them feel out of place.)}

Nhưng đời mình không thể giao cho dư luận của một nhóm quyết định thay chiều cao và chiều sâu của nó.
\textit{(But the height and depth of your life cannot be handed over to the opinion of a group.)}

Bạn bè tốt có thể góp ý cho ước mơ của em, nhưng không nên trở thành cái trần cho nó.
\textit{(Good friends may advise your dreams, but they should not become the ceiling over them.)}

\end{truyen}

\begin{baihoc}
Lắng nghe nhóm là tốt. Giao luôn tay lái ước mơ cho nhóm là nguy hiểm.
\textit{(Listening to your group is good. Handing them the steering wheel of your dream is dangerous.)}
\end{baihoc}

\begin{ghinhoanh}
\textbf{Short English to remember:}

Your friends may advise your dream, but they must not become its ceiling.

\end{ghinhoanh}

\begin{tuhocnhanh}
\textbf{Câu hỏi tự học / Self-study questions}

1. Nhóm đang tác động gì tới ước mơ của nhân vật? \textit{How is the group shaping the character's dream?}

2. Vì sao điều này nguy hiểm? \textit{Why is this dangerous?}

3. Em có đang để ai đặt trần cho giấc mơ của mình không? \textit{Are you letting anyone set the ceiling for your dream?}
\end{tuhocnhanh}
\ngancach

\section{Chúc Mừng Người Khác Mà Không Tự Nhỏ Lại}
\begin{truyen}{Chúc Mừng Người Khác Mà Không Tự Nhỏ Lại}{Celebrating Others Without Shrinking Yourself}
\chuhoa{C}{ó người trưởng thành hơn sau khi học được cách vui thật trước thành công của bạn mà không biến nó thành bằng chứng mình kém cỏi.}
\textit{(Some people grow up after learning to be genuinely happy for a friend's success without turning it into proof of their own inferiority.)}

Họ hiểu rằng mỗi thành công đẹp của người khác không cướp mất phần tương lai của mình.
\textit{(They understand that another person's good success does not steal their future.)}

Từ đó, mối quan hệ bớt căng ngầm và lòng mình cũng nhẹ hơn nhiều.
\textit{(From there, relationships become less tense and the heart grows lighter.)}

Biết chúc mừng đúng cách là dấu hiệu của một nội tâm đang khỏe dần lên.
\textit{(Knowing how to celebrate others well is a sign of an inner life becoming healthier.)}

\end{truyen}

\begin{baihoc}
Niềm vui của người khác không phải là sự mất mát của em.
\textit{(Another person's joy is not your loss.)}
\end{baihoc}

\begin{ghinhoanh}
\textbf{Short English to remember:}

Another person's joy is not your loss.

\end{ghinhoanh}

\begin{tuhocnhanh}
\textbf{Câu hỏi tự học / Self-study questions}

1. Nhân vật đã thay đổi cách nhìn nào? \textit{What new way of seeing has the character learned?}

2. Sự thay đổi đó làm lòng nhẹ ở đâu? \textit{How does that change make the heart lighter?}

3. Em có thể tập chúc mừng ai một cách thật lòng hơn? \textit{Whom can you celebrate more sincerely?}
\end{tuhocnhanh}
